\section{Introduction}
Message passing can be a general programming model, not merely an interface for
graph neural networks. Computation is organized as entities with state, relations
that determine which entities interact, local message functions, and successive
state updates. The graph need not be the input problem: it can represent the
structure of the computation itself. Tensor contractions, convolutions, particle
interactions, iterative solvers, attention and graph algorithms can all be viewed
through this lens. General message-passing models have a long history, including
the actor formalism~\cite{actors}; Tiga explores a compiler-oriented realization
of this idea for numerical and parallel programs.

The long-term goal is a common programming substrate for dense, sparse and
dynamically generated computation, rather than a collection of unrelated
operator-specific interfaces. Relations describe interaction, messages describe
local computation, reducers describe combination, and composition with control
flow describes how state evolves. A compiler can then choose how these programs
execute across kernels, memory tiers and devices while preserving their meaning.
Generality concerns what can be expressed; it does not imply that every problem
admits an efficient parallel encoding or that every encoding is supported by the
current implementation. Section~\ref{sec:general-message-passing} makes this
distinction precise.

This common form does not imply a common optimal implementation. A sparse CSR
relation needs an indexed traversal; a dense relation may need no edge array;
a generated radius relation may profit from a spatial directory. Feature width,
degree skew, gradient requirements and data residency change the suitable plan.

Tiga addresses this implementation diversity through cross-platform,
multi-backend compilation, hierarchical memory, and distributed execution.
Backend-specific lowering maps a common representation to CPU and GPU execution;
hierarchical storage decouples graph capacity from device-memory capacity; and
partition ownership with halo exchange extends execution across devices and
hosts. Memory-constrained personal systems, including laptops, motivate the
capacity goal: large graphs should not inherently require a large-memory
accelerator. The evaluation demonstrates billion-edge forward computation on
a memory-limited GPU using this capacity mechanism. Current executable backends
and extension boundaries are distinguished in Table~\ref{tab:backends}.

Tiga retains relations and reducer algebra as compiler-visible information.
It is not a model library, optimizer package or replacement for PyTorch.
Ordinary applications retain \code{torch.Tensor} and Torch autograd. Torch is
not a base installation dependency: a standalone native path supplies CPU
computation, message passing and gradients without it. Native tensors also
expose compiler/storage functionality not uniformly available through Torch.
This split is an explicit interoperability boundary, not a claim that both
interfaces support all operations identically.

Three practical costs motivate this design. First, explicitly expanding a
radius or nearest-neighbor relation can create indices, geometry and messages
that are consumed only once. Second, a graph larger than GPU memory requires
bounded storage movement even when its arithmetic is simple. Third, dividing
a graph across devices introduces halo traffic and rank imbalance; doubling
the GPU count need not reduce completion time.

The report develops a compiler representation that retains relation and reducer
semantics until traversal/kernel selection, a runtime that separates values
from residence and rank ownership, and an evaluation of the resulting execution paths.
The experiments compare time and allocation with matched Torch/PyG paths,
execute one billion explicit edges on a 16 GiB GPU, and measure the cost of
two-host NCCL graph partitioning. The results show that
generated radius traversal and exact kNN improve latency and allocation on the
measured workloads; paging extends executable capacity; and the current
equal-row two-GPU schedule is slower than the faster GPU alone. This connects
each mechanism to a measured benefit or cost rather than assuming that one
execution strategy is efficient for every relation and interconnect.

\section{Background: why graph-aware compilation?}
\label{sec:background}
\subsection{Message passing as a general computational model}
\label{sec:general-message-passing}
Message passing is broader than neighborhood aggregation in a fixed-depth
graph neural network (GNN).
In the general model, stateful entities perform local computation, communicate
values or control information, and continue through repeated transitions.
The actor formalism treats communication as an organizing principle for data,
procedures and control~\cite{actors}. Pregel gives a graph-oriented, iterative
realization through vertex programs and messages~\cite{pregel}. Tiga uses
structured relations and reducers rather than implementing arbitrary actor
mailboxes or adopting Pregel's execution protocol.
In graph learning, Message Passing Neural Networks (MPNNs) formalize learned
messages and aggregation across repeated node updates~\cite{mpnn}. Tiga uses
this local-computation interface as a compiler input, while also supporting
non-neural message functions such as stencil updates and weighted sums.

\paragraph{Universality and its conditions.}
With sufficiently expressive local transitions, repeatable control and
arbitrarily extensible state, a message-passing model can in principle express
any computable algorithm. A constructive way to see this is to encode a Turing
machine's tape as a chain of entities. Each entity stores a tape symbol; a
message carries the head's finite control state. Receiving the head message
updates the symbol and sends the next head state to the left or right neighbor.
Allowing the chain and the number of transitions to grow reproduces the machine's
computation. This is an expressiveness argument for the general model, not a
universality proof for Tiga's present bounded loops, fixed-shape tensors or
finite-memory backends. It refers to computable problems, not undecidable ones,
and says nothing by itself about time, memory or useful parallelism.

\paragraph{A common language for numerical programs.}
The practical opportunity is to retain a useful computational structure without
requiring every application to start from an explicit graph dataset. Elementwise
maps use independent entities or self-relations. Matrix multiplication connects
each output row to contributing source rows, with weighted vector messages and
a sum reducer. Convolution and finite-difference stencils use regular local
relations; particles use coordinate-dependent radius or kNN relations; attention
uses dense or masked relations and a normalization-aware reducer. Repeated
message passing with state updates expresses time stepping, iterative solvers
and graph traversals. These are representations within a shared conceptual model,
not a claim that every specialized operation has an equally optimized Tiga path.
The language ambition is to make such programs composable and their interactions
compiler-visible; the implementation and evaluation below establish concrete
slices of that ambition.

\paragraph{Why a compiler is needed.}
An edge function can be only a multiply, yet its execution may require reading
an index, gathering a feature vector, creating an edge-sized temporary and
scattering into a destination. For $E$ edges and feature width $F$, explicitly
materializing messages requires storage proportional to $EF$, even when the
final output has only $NF$ values. A fused traversal can instead accumulate
messages directly. Fusion is not universally optimal: a library sparse matrix
operation, feature tiling, recomputation or a different graph layout may win
depending on shape, degree distribution and reuse. Consequently the relevant
compiler question is which semantic information remains available when each
decision is made, not whether a system uses the phrase ``message passing.''

PyTorch Geometric supplies a graph-learning programming model with gather--scatter
and optimized sparse execution paths~\cite{pyg,pyg2}. Graphiler makes the compiler connection explicit:
its message-passing data-flow graph annotates graph-related tensor residency and
movement, enabling reductions in redundant computation and intermediate
storage~\cite{graphiler}. Tiga likewise retains source-, edge- and
destination-associated semantics, and carries these roles into relation traversal
and differentiation.

\subsection{Graph and feature schedules, fusion, and training}
FeatGraph composes sparse traversal templates with feature-dimension UDFs and
schedules~\cite{featgraph}. It shows why graph traversal and within-feature
computation should be optimized together. Seastar offers vertex-centric Python
programming and generates fused forward and backward GPU kernels~\cite{seastar}.
Tiga shares the goal of fusing local computations, and combines it with
explicit relation constructors and storage/ownership representations.

GALA is especially relevant to a multi-level design. It separates data and
compute representations, composes intra-operator and inter-operator
transformations, and reasons about training-specific optimization opportunities
across forward and backward computations~\cite{gala}. Tiga separates relation,
traversal and execution representations; its current evaluation focuses on
forward workloads, with path-specific gradient implementations described in
the differentiation section. Whole-training-program optimization remains beyond
that evaluated scope.

GraphIt separates the graph algorithm from scheduling choices~\cite{graphit}.
That distinction helps explain why a relation-level operation should not
prematurely encode threads, blocks or traversal order. Tiga additionally seeks
to preserve numerical field operations, reducer state and differentiation
requirements alongside traversal scheduling.

\subsection{Relations and sparse iteration beyond GNNs}
Ebb represents geometric domains through a relational data model and separates
simulation code, domain libraries and parallel runtime implementation~\cite{ebb}.
Simit connects hypergraph-local computation to global tensor algebra for physical
simulation~\cite{simit}. They are important precedents for a graph language that
is not limited to neural networks. Tiga's particle, solver and stencil examples
belong in this lineage of relational numerical programming.

TACO derives compound dense/sparse tensor kernels from high-level algebra using
iteration and merge structures~\cite{taco}. Its format-abstraction work further
separates logical computation from sparse storage levels~\cite{sparseformats}.
SparseTIR develops composable sparse formats and transformations for deep-learning
operators~\cite{sparsetir}. These works motivate keeping coordinate traversal
distinct from physical layout and execution mapping. Tiga's Iter IR preserves
relation-coordinate structure for its supported traversal families; the
implementation does not provide a general sparse-format algebra.

Triton supplies a lower-level tiled compilation interface~\cite{triton}, while
MLIR supplies infrastructure for composing dialects and lowering passes~\cite{mlir}.
Tiga uses both: graph-aware representations are upstream of a current NVIDIA
provider. The provider performs device-code compilation after Tiga's graph-aware
transformations.

\subsection{Capacity: running the graph at all}
Reducing execution latency and reducing the resident working set are related but
distinct goals. If a graph does not fit a resource budget, a faster in-memory
kernel is insufficient. GraphChi established disk-based large-graph computation
on a single machine through partitioned processing and sliding
windows~\cite{graphchi}. X-Stream uses edge-centric streaming partitions to favor
sequential access over random graph-data access~\cite{xstream}. These systems are
relevant even though they do not share Tiga's differentiable Tensor interface.

Legion separates logical regions and task privileges from placement and physical
instances, providing a model for reasoning about locality, independence and
data movement~\cite{legion}. Tiga's storage versions, task dependencies and halo
ownership apply this separation to the graph runtime described in Section~\ref{sec:runtime}.

These systems distinguish executing faster from making execution possible.
Tiga evaluates both, plus the cost of distributing the same graph. The capacity
experiment traverses explicit disk-backed CSR through host staging to CUDA;
the distributed experiment measures rank ownership and halo movement separately.
Neither a transfer-bandwidth microbenchmark nor a successful communicator setup
alone answers these application-level questions.

FlashAttention is a related example of IO-aware exact
computation~\cite{flashattention}: avoiding an intermediate can matter as much as
reducing arithmetic. Tiga's dense relations and streaming reducers expose
similar scheduling opportunities, while optional tile pruning changes the
mathematical contract and is distinct from exact attention. The experiments in
this report concern graph aggregation and do not claim attention-library parity.

\begin{table}[ht]
\centering\small
\begin{tabularx}{\linewidth}{l Y Y}
\toprule
Prior work & Relevant mechanism & Connection to Tiga \\
\midrule
Graphiler~\cite{graphiler} & Graph-aware dataflow/residency & Source, edge and destination field roles \\
FeatGraph~\cite{featgraph} & Sparse traversal + feature schedules & Joint graph/feature traversal selection \\
Seastar~\cite{seastar} & Vertex-centric fused forward/backward & Message fusion and derivative generation \\
GALA~\cite{gala} & Composed data/compute and training transforms & Layered representation and transformation scope \\
Ebb / Simit~\cite{ebb,simit} & Relational simulation / hypergraph algebra & Numerical programs beyond GNNs \\
GraphIt~\cite{graphit} & Algorithm/schedule separation & Separate meaning from execution mapping \\
TACO / SparseTIR~\cite{taco,sparsetir} & Sparse iteration, formats, transformations & Relation traversal separate from storage \\
GraphChi / X-Stream~\cite{graphchi,xstream} & Disk-oriented partitioned execution & Bounded graph pages and storage traffic \\
Legion~\cite{legion} & Logical regions and physical placement & Storage instances, ownership and dependencies \\
\bottomrule
\end{tabularx}
\caption{Prior work and its relationship to Tiga's design.}
\end{table}
\FloatBarrier


\section{Programming model}
\label{sec:programming-model}
\subsection{Relations, fields and messages}
Let source entities and destination entities be distinct finite sets. A relation
is an ordered collection of edges with source and destination maps. Duplicate
edges remain separate messages. For a destination, the computation is
$$
 m_e = f(x_{s(e)}, z_{d(e)}, a_e;\theta),\qquad
 y_i = g\!\left(z_i,\operatorname{finalize}\left(
       \bigoplus_{e:d(e)=i}\operatorname{lift}(m_e)\right);\theta\right).
$$
The reducer supplies its identity, lift, combine and finalize behavior.
An empty row starts from the identity; finalization determines the resulting
empty-row convention. This convention is reducer-specific. Declaring a reducer
does not define which edge-local values enter it: the edge function does.
Simple sum messages return a tensor directly. A structured streaming reducer
such as online softmax currently packages a score and value using an explicit
reducer call inside the edge function. That call does not execute an aggregation
once per edge; it binds the message inputs for capture.

CSR is destination-oriented: row boundaries enumerate destinations and column
indices identify sources. The example below has two edges into destination 0,
one into destination 1, and none into destination 2. A self edge is ordinary;
there is no automatic reverse edge.
\noindent\begin{minipage}{\linewidth}
\begin{lstlisting}[language=Python]
import torch
import tiga as tg

class WeightedSum(tg.MessagePassing):
    reducer = tg.sum()
    def edge(self, src, dst, edge):
        return src.x * edge.w

graph = tg.Graph.from_csr(
    torch.tensor([0, 2, 3, 3]), torch.tensor([0, 1, 1]),
    num_src=2, validate="full")
x = torch.tensor([2., 3.], requires_grad=True)
w = torch.tensor([4., 5., 2.], requires_grad=True)
y = WeightedSum()(graph=graph, src={"x": x}, dst={}, edge={"w": w})
dx, dw = torch.autograd.grad(y.sum(), (x, w))
# y=[23,6,0], dx=[4,7], dw=[2,3,3]
\end{lstlisting}
\end{minipage}

\begin{figure}[ht]
\centering
\begin{tikzpicture}[>=Stealth, every node/.style={font=\small}]
\node[draw,circle] (s0) at (0,1.0) {source 0};
\node[draw,circle] (s1) at (0,-1.0) {source 1};
\node[draw,circle,very thick,draw=tigateal] (d0) at (5,0) {destination 0};
\draw[->,thick,color=tigateal] (s0) -- node[above,sloped] {$2\times4=8$} (d0);
\draw[->,thick,color=tigateal] (s1) -- node[below,sloped] {$3\times5=15$} (d0);
\node[right=0.5cm of d0,text=tigateal] {$8+15=23$};
\end{tikzpicture}
\caption{One destination row: two edge messages become one result. Other
destinations are intentionally omitted.}
\end{figure}

\subsection{Composition and topology}
Dense and triangular constructors describe relation structure, while
\code{Graph.cat} forms disjoint blocks by offsetting source and destination
IDs independently. For example, concatenating triangular graphs of sizes two
and three represents independent causal sequences. It introduces no cross-block
edges and currently materializes CSR; generic composition is not yet an
implicit-storage optimization. Cumulative sequence boundaries remain a
compatibility constructor and must start at zero.

Radius and nearest-neighbor builders depend on coordinates. Their relation
realization may differ by execution path. Topology selection is discrete;
the selected graph is a fixed snapshot for derivative semantics. In-place
mutation of captured indices is not a supported topology-update mechanism.
Transpose and composition can reorder edges, so explicit edge fields require
matching reordering. A constructor accepting an input does not certify every
consumer, provider or gradient combination.
